% ============================================================
% Section 7 — Correctness Properties
% ============================================================
\section{Correctness Properties}
\label{sec:rs-correctness}

This section establishes the formal correctness guarantees of the Recursive Spawning Protocol with respect to its core objective: the elimination of autonomous drift. We present three theorems covering drift prevention, chain integrity, and termination, together with the structural invariants that underpin each proof. The proofs are proof sketches; full formalization is a matter for future machine-verified proof work.

% ----------------------------------------------------------
\subsection{Drift Prevention}
\label{sec:rs-drift-prevention}

\begin{thrm}[Drift Prevention]
\label{thrm:drift-prevention}
Let $C = \langle n_0, n_1, \ldots, n_k \rangle$ be a Recursive Spawning chain rooted at $n_0$ with human anchor $h(n_0) = h(n_k)$. Then for all $i$ ($0 \leq i \leq k$), the operating scope $R(n_i)$ is a descendant refinement of the intent of $h(n_0)$. Equivalently: no node in $C$ can exhibit behavior outside the authorized intent of its human anchor.
\end{thrm}

\begin{proof}[Proof Sketch]
The proof proceeds by structural induction on chain depth, where each inductive step is enforced at runtime by the Fact Layer pre-commit hook.

\medskip
\noindent\textit{Base case ($i=0$).} $R(n_0)$ is defined by $h(n_0)$ at provisioning. By the Purpose Layer's definition of authorized operating scope, $R(n_0)$ is a descendant refinement of $h(n_0)$'s intent by construction. No drift exists at the root.

\medskip
\noindent\textit{Inductive step.} Assume $R(n_i)$ is a descendant refinement of $h(n_0)$'s intent for some $i$ with $0 \leq i < k$. For $n_i$ to spawn $n_{i+1}$, the Fact Layer pre-commit hook verifies two conditions before the spawn is recorded:

\begin{enumerate}
    \item[(a)] $R(n_{i+1}) \subseteq R(n_i)$ --- scope refinement. The Purpose Layer's spawn authorization policy $P_{\mathrm{spawn}}$ requires strict scope subset; there is no authorized provision for lateral or upward scope expansion.
    \item[(b)] $\mathrm{depth}(n_{i+1}) \leq D_{\mathrm{max}}$ --- depth constraint, also enforced by the pre-commit hook.
\end{enumerate}

The Fact Layer rejects any spawn operation violating (a) or (b). No actor, including the Bounds Engine, can execute a rejected spawn. Therefore $R(n_{i+1}) \subseteq R(n_i)$ holds for the child by structural enforcement, not by policy convention.

By the inductive hypothesis, $R(n_i)$ is a descendant refinement of $h(n_0)$'s intent. Since $R(n_{i+1}) \subseteq R(n_i)$, transitivity of subset gives that $R(n_{i+1})$ is also a descendant refinement of $h(n_0)$'s intent. Hence $n_{i+1}$ does not drift. By induction, no node in $C$ drifts. \qed
\end{proof}

\medskip
The drift prevention guarantee does not rest on the goodness or reliability of any actor within the system. It rests on the conjunction of three architectural facts, each enforced by a distinct BX3 layer:

\begin{enumerate}
    \item[(1)] \textbf{Immutable parent pointer} (Fact Layer). $\alpha(n)$ is established once at provisioning and is thereafter immutable under the Fact Layer write protocol. Scope inheritance is traceable to a unique authorized lineage, not a manipulable graph. This eliminates the re-parenting drift pathway (Failure Mode 3, Section~\ref{sec:drift-failure-modes}).

    \item[(2)] \textbf{Forensic ledger} (Fact Layer). Every scope assignment is recorded in the append-only, hash-chained ledger. The ledger makes each step of the refinement chain an empirically verifiable, tamper-evident record. No action can occur without being logged; no log entry can be altered after the fact.

    \item[(3)] \textbf{Chain terminates at Purpose Layer} (Purpose Layer). The exclusive human terminus closes the inductive argument at a finite bound. There is no machine-actor alternative terminus; $h(n_i) = h(n_0)$ for all $i$. This prevents the chain from self-authorizing its own continuation beyond the human anchor's intent.
\end{enumerate}

With all three conjuncts present, drift is architecturally impossible --- not statistically unlikely, not conventionally prevented. If any conjunct is absent, the proof structure collapses. The depth-limit insufficiency theorem (Section~\ref{sec:drift-depth-limits}) establishes that condition (3) alone is insufficient without (1) and (2); the three conditions are jointly necessary and jointly sufficient.

% ----------------------------------------------------------
\subsection{Chain Integrity}
\label{sec:rs-chain-integrity}

\begin{thrm}[Chain Integrity]
\label{thrm:chain-integrity}
For any Recursive Spawning chain $C = \langle n_0, n_1, \ldots, n_k \rangle$, the parent pointer sequence $\alpha(n_i) = n_{i-1}$ holds for all $1 \leq i \leq k$, and no node configuration in $C$ has been modified after provisioning.
\end{thrm}

\begin{proof}[Proof Sketch]
The proof has two components.

\medskip
\noindent\textit{Topological integrity.} Each spawn event $n_{i-1} \xrightarrow{\mathrm{spawn}} n_i$ establishes $\alpha(n_i) = n_{i-1}$ as an immutable Fact Layer property at provisioning time. The Fact Layer write protocol rejects any operation attempting to modify $\alpha(n_i)$ after provisioning. Since the chain is constructed exclusively through authorized spawn events and no parent pointer is modified after establishment, $\alpha(n_i) = n_{i-1}$ holds for all $i \geq 1$ by construction. The parent pointer is not merely access-controlled; it is structurally immutable. There is no privileged actor, including the Purpose Layer after provisioning, capable of overwriting $\alpha(n)$. This eliminates the re-parenting drift pathway by making retroactive chain rewriting structurally impossible.

\medskip
\noindent\textit{Node immutability.} Each node $n_i$ is provisioned with a fixed configuration $(R(n_i), P_{\mathrm{spawn}}(n_i), h(n_i))$ established by the Purpose Layer and recorded in the Fact Layer authorization ledger. The Fact Layer enforces immutability of these parameters after provisioning; the write protocol rejects any modification request from any actor. The Purpose Layer may issue a \texttt{RESOLVE} directive that redirects a node's behavior, but this directive is recorded as a new ledger entry --- a state transition --- not as an alteration of the provisioned configuration.

The forensic ledger's hash-chaining provides evidentially for chain integrity: any attempt to insert a falsified spawn event or reorder existing entries breaks hash continuity and is detectable by any observer with ledger read access. Chain integrity is not merely a property of the current chain state; it is a property of the entire historical record, enforceable at any future time. \qed
\end{proof}

\begin{prop}[Non-Repudiation of Spawn]
\label{prop:non-repudiation}
A spawn event recorded in the Fact Layer of the parent node constitutes irrefutable evidence of the child's existence, its spawn parameters, and its delegation chain at the time of spawn.
\end{prop}

\begin{proof}
The spawn event is recorded in the parent's Fact Layer immediately upon child creation, before the child begins execution. The record is immutable once written. The parent cannot deny having created the child, because the Fact Layer entry is timestamped, hash-chained, and signed by the parent's own Fact Layer key material. The child cannot deny its chain of delegation, because its Purpose Layer chain anchor matches the parent's Fact Layer spawn record. \qed
\end{prop}

% ----------------------------------------------------------
\subsection{Termination}
\label{sec:rs-termination}

\begin{thrm}[Termination]
\label{thrm:termination}
Every Recursive Spawning chain $C = \langle n_0, n_1, \ldots, n_k \rangle$ terminates --- either at a resolution state or at the escalation threshold --- within a finite number of spawn steps bounded by $\max(D_{\mathrm{max}}, D_{\mathrm{ESCALATE}})$.
\end{thrm}

\begin{proof}[Proof Sketch]
Two lemmas combine to establish the theorem.

\medskip
\noindent\textit{Lemma 1 (Depth-bounded growth).} The chain grows only through spawn events. By the Bounds Engine depth enforcement and the Fact Layer pre-commit hook, any spawn event producing $|C| \geq D_{\mathrm{max}}$ is rejected. Therefore the chain cannot exceed depth $D_{\mathrm{max}} - 1$ through autonomous spawn. Rejected spawns are logged; the Bounds Engine transitions to \texttt{ESCALATING}. Autonomous chain growth is thus bounded by $D_{\mathrm{max}}$.

\medskip
\noindent\textit{Lemma 2 (Escalation-triggered halt).} If the chain reaches $D_{\mathrm{ESCALATE}}$ before resolving the exception condition, the protocol enters \texttt{ESCALATING} state and suspends all autonomous spawn activity. The Bounds Engine awaits a directive from $h(n_k)$ --- the human anchor. No further spawn events occur until $h(n_k)$ issues \texttt{ACK}, \texttt{RESOLVE}, or \texttt{REJECT}. Since $h(n_k) = h(n_0)$ is a human principal, the directive arrives in finite time. \texttt{REJECT} terminates the chain; \texttt{RESOLVE} redirects behavior and records a new ledger entry; \texttt{ACK} permits continued operation under continued Purpose Layer authorization.

\medskip
\noindent\textit{Theorem conclusion.} The chain grows autonomously at most $D_{\mathrm{max}} - 1$ spawn steps. If it reaches $D_{\mathrm{ESCALATE}}$ before resolution, it halts and escalates. In all execution paths, the chain reaches a terminal state within $\max(D_{\mathrm{max}}, D_{\mathrm{ESCALATE}})$ spawn steps. Since both $D_{\mathrm{max}}$ and $D_{\mathrm{ESCALATE}}$ are finite integers defined at provisioning, the chain terminates in finite time. \qed
\end{proof}

\begin{lemm}[Orphan Safety]
\label{lemm:orphan-safety}
No node can be orphaned without its parent entering a terminal state.
\end{lemm}

\begin{proof}
A node becomes orphaned when its parent enters terminal state without issuing a termination event. Under RSP, the transition to terminal state is defined as a structured procedure (Theorem~\ref{thrm:termination}) that requires the issuance of termination events to all active children as its first step. This procedure is enforced by the Bounds Engine: a node cannot enter terminal state without the Bounds Engine confirming that all active children have received termination events. A node whose parent has terminated without issuing such events would have a Bounds Engine violation recorded in its Fact Layer indicating a parent-termination anomaly, detectable by any external auditor. \qed
\end{proof}

\begin{cor}[No Ghost Authority]
\label{cor:ghost-authority}
No node can be spawned without an authorized delegation chain terminating at a human principal's Purpose Layer.
\end{cor}

\begin{proof}
Directly from Theorem~\ref{thrm:drift-prevention} and Proposition~\ref{prop:non-repudiation}. If a node could be spawned without an authorized delegation chain, it would be a drifted node by definition, contradicting Theorem~\ref{thrm:drift-prevention}. \qed
\end{cor}

% ----------------------------------------------------------
\subsection{Collective Guarantee}
\label{sec:rs-collective-guarantee}

The three correctness properties are mutually reinforcing and jointly enforced by the BX3 Framework's three-layer architecture. Drift prevention ensures that scope refinement is the only authorized direction of chain growth: each child scope is a strict subset of its parent, traceable to the human anchor's intent through the forensic ledger. Chain integrity ensures that chain topology is stable and tamper-evident: the parent pointer sequence is immutable, node configurations are unmodifiable after provisioning, and hash-chaining makes any tampering detectable. Termination ensures that the chain cannot grow indefinitely and cannot sustain an infinite spawn-escalate cycle: the depth bound is enforced structurally, and the escalation path terminates in human judgment.

Together these properties guarantee that a Recursive Spawning chain is an accountable, bounded, and auditable refinement process: it grows only in authorized directions (drift prevention), preserves the topology established at provisioning (chain integrity), halts within a Purpose Layer-defined depth bound (termination), and maintains a complete forensic record of every decision point (Fact Layer). This is the operational expression of the upstream BX3 accountability guarantee applied to recursive agent population dynamics.
